% ============================================================
\subsection{Ad Serving and Matching}
\label{subsec:ad-serving-matching}
% ============================================================

This subsection describes the path a user turn takes from the moment the
orchestrator decides that an advertisement may be appropriate, to the moment a
small ranked set of candidate ads is handed back to the response synthesis
stage. We treat the problem as a two-stage retrieval-and-ranking pipeline: a
high-recall semantic retrieval over a vector index, followed by a
business-rules re-ranking step that enforces the economic and policy
constraints of the platform. This decomposition mirrors the structure of
modern web search and recommendation systems, where embedding-based recall is
combined with a lightweight learned or hand-tuned re-ranker
\cite{karpukhin2020dpr}, but is adapted here to the constraints of a
multi-turn conversational setting.

\subsubsection{Ad Inventory, Campaigns, and Advertisers (PostgreSQL + pgvector)}
\label{subsubsec:ad-inventory-store}

The ad inventory is the system of record for every entity an advertiser cares
about: the advertiser account itself, the campaigns it funds, and the
individual creatives that may be served. We store all three in a single
PostgreSQL instance, extended with the \texttt{pgvector} extension to support
approximate nearest-neighbour search over dense embeddings. The decision to
co-locate relational and vector data in PostgreSQL, rather than to introduce a
dedicated vector database such as Pinecone or Weaviate, is deliberate. At the
scale we target in this report (on the order of $10^3$--$10^5$ ads),
\texttt{pgvector} with an IVFFlat index comfortably meets our latency budget,
and keeping ads, budgets, and embeddings in the same transactional store
removes an entire class of consistency bugs that arise when budget
decrements and embedding lookups live in different systems.

The schema is organised around three tables. The \texttt{advertisers} table
holds account-level metadata and contact information. The \texttt{campaigns}
table groups ads under a shared budget envelope, bid strategy, and flight
window, and references an advertiser via a foreign key. The \texttt{ads}
table holds the per-creative fields that the matching pipeline actually reads
at request time:

\begin{itemize}[leftmargin=1.4em,itemsep=2pt]
    \item \textbf{Identity and content:} \texttt{ad\_id}, \texttt{title},
          \texttt{description}, \texttt{cta\_url}, \texttt{product\_category}.
    \item \textbf{Targeting:} a free-text \texttt{targeting\_context} field
          that captures the kinds of conversational situations the advertiser
          wants to reach (e.g.\ ``users researching marathon training shoes'').
    \item \textbf{Economics:} \texttt{bid\_cpm}, \texttt{remaining\_budget},
          \texttt{daily\_cap}, and \texttt{active} flag.
    \item \textbf{Embedding:} a \texttt{vector(1536)} column populated at
          ingestion time by the embedding pipeline of
          Section~\ref{subsubsec:embedding-pipeline}.
\end{itemize}

An IVFFlat index is built on the embedding column using cosine distance as
the similarity operator. The choice of cosine over inner product is motivated
by the fact that the OpenAI \texttt{text-embedding-3-small} embeddings used in
this work are not unit-normalised in a way that makes raw inner products
meaningful for ranking across heterogeneous ad copy lengths. Cosine
similarity is invariant to embedding magnitude and therefore yields more
stable rankings when ad descriptions vary in length, which is the common case
in our inventory.

The \texttt{events} table, although logically part of the feedback loop
described in Section~\ref{subsec:feedback-loop}, is co-located in the same
database so that impression and click writes can participate in the same
transaction as the budget decrement on the corresponding ad row. This is
what allows us to make a strong claim later: a charged impression and a
budget deduction either both happen or neither does.

\subsubsection{Embedding and Semantic Matching Pipeline}
\label{subsubsec:embedding-pipeline}

Semantic matching in the system is built on the principle that the ad and
the conversational context should be projected into the \emph{same} vector
space, but from \emph{different} textual surfaces. On the ad side, we embed
a concatenation of the title, description, product category, and targeting
context; on the query side, we embed a structured serialisation of the
output of the context extraction module
(Section~\ref{subsubsec:context-extraction}), not the raw user message.

This asymmetry is important and worth justifying. The raw user message is a
poor query for two reasons. First, it is often short, casual, and lacks the
domain vocabulary that ad copy is written in --- a user who says ``my knees
are killing me after long runs'' shares almost no surface tokens with an ad
titled ``Nike Pegasus 41 --- Cushioned Daily Trainer''. Second, intent in a
conversational setting is distributed across turns, so any embedding of a
single message systematically discards the accumulated context that the
session store has been maintaining. By embedding the structured context
object instead --- which contains the extracted intent, topics, entities, and
purchase signal --- we obtain a query representation that is both richer and
more lexically aligned with the way advertisers describe their products.

Formally, let $f_\theta : \Sigma^* \to \mathbb{R}^{1536}$ denote the
embedding model. For an ad $a$ with textual fields $T(a)$, the ingestion-time
embedding is $\mathbf{e}_a = f_\theta(T(a))$, computed once and persisted in
the \texttt{ads} table. For a user turn $m_t$ with conversation history $H$,
let $C = \texttt{extract\_context}(m_t, H)$ be the structured context, and
let $S(C)$ be its canonical serialisation. The query embedding is then
$\mathbf{q} = f_\theta(S(C))$, computed at request time. Retrieval reduces to
finding the top-$N$ ads under cosine similarity:
\[
    \mathit{Retrieve}(\mathbf{q}, N) \;=\;
    \mathop{\mathrm{arg\,top}\text{-}N}_{a \in \mathcal{A}}\;
    \frac{\mathbf{q}^\top \mathbf{e}_a}{\lVert \mathbf{q}\rVert\,\lVert \mathbf{e}_a\rVert}.
\]
We set $N = 20$, which empirically gives the downstream re-ranker enough
diversity to enforce budget and frequency constraints without exhausting the
inventory in low-fill regimes.

Two practical optimisations are worth noting. First, the embedding call for
the query is the second most expensive synchronous step in the pipeline
(after response generation), so the orchestrator issues it in parallel with
the tail end of context extraction whenever the structured context becomes
incrementally available. Second, the IVFFlat index is warmed at process
start by issuing a dummy nearest-neighbour query, which avoids a cold-cache
penalty on the first real user turn after a deployment.

\subsubsection{Re-Ranking, Budget Gating, and Business Logic}
\label{subsubsec:reranking-business-logic}

The output of semantic retrieval is a high-recall set of $N=20$ candidates
that are \emph{topically} plausible but not yet economically or contextually
admissible. The re-ranker is the component that turns this candidate set into
a small ranked list ready for the response synthesis stage. It performs three
jobs in sequence: \emph{filtering}, \emph{scoring}, and \emph{selection}.

\paragraph{Filtering.} A candidate ad $a$ is admissible only if it satisfies
the conjunction of the following predicates: (i) $a.\texttt{active}$ is true,
(ii) $a.\texttt{remaining\_budget} > a.\texttt{bid\_cpm}/1000$ (i.e.\ the
campaign can afford one more impression), and (iii) $a.\texttt{ad\_id}$ does
not appear in the \texttt{ads\_shown\_this\_session} list maintained by Redis.
The third predicate is the per-session frequency cap and is the simplest but
most user-visible safeguard against repetition fatigue: a user who has
already seen a Nike Pegasus mention in turn 3 will not see it again in turn 5,
even if it remains the highest-similarity candidate. Filtering is performed
in application code rather than as part of the SQL query because the
session-level shown set lives in Redis, not in PostgreSQL, and we prefer to
keep that boundary explicit.

\paragraph{Scoring.} For each surviving candidate we compute a composite
score that blends semantic relevance with the platform's economic objective
and the inferred user state. Let $\mathrm{sim}(a)$ denote the cosine
similarity between $\mathbf{q}$ and $\mathbf{e}_a$, and let
$b_{\max} = \max_{a' \in \mathit{candidates}} a'.\texttt{bid\_cpm}$ be the
maximum bid in the surviving set. We define
\[
    \mathit{score}(a) \;=\;
    0.5 \cdot \mathrm{sim}(a)
    \;+\; 0.3 \cdot \frac{a.\texttt{bid\_cpm}}{b_{\max}}
    \;+\; 0.2 \cdot \mathit{purchase\_signal}(C),
\]
where $\mathit{purchase\_signal}(C) \in [0,1]$ is the scalar field produced
by the context extraction module. Normalising bid by $b_{\max}$ rather than
by a global constant keeps the bid term on the same scale as the other two
across very different inventories, which is important during early-stage
testing when the inventory is small and bid magnitudes are arbitrary.

The weights $(0.5, 0.3, 0.2)$ are deliberately chosen to make relevance the
dominant factor: even the highest bidder cannot displace a much more relevant
ad, and a high purchase signal can act as a tie-breaker between two ads of
similar relevance. We treat these weights as hand-tuned hyperparameters in
this report; learning them from logged engagement data is a natural
extension and is discussed in Section~\ref{subsec:future-work}.

\paragraph{Selection.} The top $k=3$ scored candidates are returned to the
orchestrator. Returning more than one allows the response synthesis prompt
(Section~\ref{subsubsec:response-synthesis}) to gracefully fall back to a
second-choice ad if the LLM judges the top candidate to be a poor
contextual fit, without requiring a second round-trip to the database.

Algorithm~\ref{alg:reranking} summarises the complete procedure as
implemented in \texttt{get\_candidate\_ads}, which is the function exposed to
the AI workstream through the interface contract of
Table~\ref{tab:interface-contracts}.

